% \begin{document}
\chapter{Attacchi di Buffer Overflow}

\section{Introduzione}

\subsection{Buffer overflow attacks}

Il \textbf{software} \`e noto per avere \textbf{bug}:
\begin{itemize}
  \item Funzionalit\`a che non vanno come previsto o non funzionano affatto.
  \item \textbf{Crash} che causano inaffidabilit\`a e perdita di dati.
\end{itemize}

\nt{
  Per un attacker, \textbf{i bug nel software sono delle opportunit\`a}.
}

\subsection{Exploits}

\dfn{Exploit}{
  \textbf{Bug del software trasformati in armi}: l'utilizzo di errori di programmazione a vantaggio di un utente malintenzionato.
}

Usi tipici degli exploit:
\begin{itemize}
  \item Aggirare i controlli di \textbf{autenticazione} e \textbf{autorizzazione}
  \item \textbf{Elevare i privilegi} (ad admin o root)
  \item Dirottare i programmi per eseguire \textbf{codice non intenzionale / arbitrario / shellcode}
  \item Consentire l'\textbf{accesso persistente e non autorizzato} ai sistemi
\end{itemize}

\subsection{Storia degli incidenti pi\`u importanti}

\begin{center}
\begin{tabular}{|l|p{10cm}|}
\hline
\textbf{Anno} & \textbf{Evento} \\
\hline
1988 & Il Morris Internet Worm us\`o un exploit di buffer overflow in \texttt{fingerd}. \\
1995 & Buffer overflow in NCSA httpd 1.3 pubblicato sulla mailing list Bugtraq da Thomas Lopatic. \\
1996 & Aleph One pubblica \textit{Smashing the Stack for Fun and Profit} su \textit{Phrack} \#49. \\
2001 & Worm \textbf{Code Red}: buffer overflow in Microsoft IIS 5.0. \\
2003 & Worm \textbf{Slammer}: buffer overflow in Microsoft SQL Server 2000. \\
2004 & Worm \textbf{Sasser}: buffer overflow nel LSASS di Windows 2000/XP. \\
\hline
\end{tabular}
\end{center}

\section{Definizione di buffer overflow}

\dfn{Buffer Overflow (NISTIR 7298, 2019)}{
  Una condizione su un'interfaccia in cui \`e possibile inserire \textbf{pi\`u input in un buffer} (o area di conservazione dei dati) rispetto alla \textbf{capacit\`a assegnata}, sovrascrivendo altre informazioni.

  Gli avversari sfruttano tale condizione per mandare in crash un sistema o per inserire codice appositamente predisposto che consente loro di ottenere il controllo del sistema.
}

\nt{
  Detto anche \textbf{buffer overrun} o \textbf{buffer overwrite}.
}

\paragraph{Esempio --- Buffer da 8 byte}: buffer che dovrebbe contenere \texttt{USERNAME}, ma riceve \texttt{USERNAME12}:
\begin{verbatim}
| Buffer (8 bytes)           | Overflow |
| U | S | E | R | N | A | M | E | 1 | 2 |
  0   1   2   3   4   5   6   7   8   9
\end{verbatim}
I byte \texttt{1} e \texttt{2} finiscono \textbf{oltre il buffer}, sovrascrivendo memoria adiacente.

\section{Esecuzione di un programma}

\subsection{I compilatori}

\[
  [\texttt{C}] \to \boxed{\text{Front End}} \to \boxed{\text{IR}} \to \boxed{\text{Back End}} \to [\texttt{ASM/BIN}]
\]

\begin{itemize}
  \item I computer \textbf{non eseguono codice sorgente} (Java, C++, ...).
  \item Eseguono \textbf{codice macchina} (machine code).
  \item I \textbf{compilatori} traducono il codice sorgente in codice macchina.
  \item L'\textbf{assembly} \`e codice macchina leggibile dall'uomo.
\end{itemize}

\subsection{Esempio C $\to$ x86-64}

\begin{verbatim}
#include <stdio.h>
int main(int argc, char** argv) {
    int i;
    if (argc > 1) {
        for (i = 1; i < argc; ++i) {
            puts(argv[i]);
        }
    }
    else {
        puts("Hello world");
    }
    return 1;
}
\end{verbatim}

In assembly x86-64 (estratto):
\begin{verbatim}
push   rbp
mov    rbp,rsp
sub    rsp,0x20
mov    DWORD PTR [rbp-0x14],edi
mov    QWORD PTR [rbp-0x20],rsi
cmp    DWORD PTR [rbp-0x14],0x1
jle    400578 <main+0x4b>
...
call   400410 <puts@plt>
...
leave
ret
\end{verbatim}

\subsection{Memoria del computer}

I \textbf{programmi in esecuzione} sono presenti nella memoria. Layout (Virtual Memory $\sim 4$ GB):

\begin{verbatim}
4 GB +----------------------+
     |  Operating System    |
     +----------------------+
     |   Data Memory        |
     |   (Variables)        |
     +----------------------+
     |       ...            |
     +----------------------+
     |   Program Memory     |
     |   (Code)             |
   0 +----------------------+
\end{verbatim}

\begin{itemize}
  \item \textbf{Program Memory}: il codice del programma.
  \item \textbf{Data Memory}: variabili, costanti e altre cose necessarie.
  \item \textbf{OS Memory}: sempre disponibile per le \textbf{system calls} (aprire un file, eseguire un altro programma, stampare, ecc.).
\end{itemize}

\subsection{Memoria del programma e Instruction Pointer (IP)}

\dfn{Instruction Pointer (IP / PC)}{
  Registro della CPU che tiene traccia dell'istruzione corrente. Detto anche \textit{Program Counter}.
}

\begin{verbatim}
0: integer count(string s, character c) {
     integer counter=0;
     integer pos;
1:   for (pos = 0; pos < length(s); pos = pos + 1) {
2:     if (s[pos] == c) counter = counter + 1;
3:   }
4:   return counter;
5: }

6: void main(integer argc, strings argv) {
7:   count("testing", "t"); // should return 2
8: }
\end{verbatim}

L'IP punta inizialmente alla riga 7 (chiamata di \texttt{count}). Il codice sta in Program Memory; le variabili (\texttt{counter}, \texttt{pos}) e le costanti (\texttt{"testing"}, \texttt{"t"}) sono salvate nella Data Memory.

\section{Lo Stack e la chiamata di funzione}

\subsection{Lo Stack}

La \textbf{memoria dati} \`e organizzata usando una struttura specifica: lo \textbf{stack}.

\begin{verbatim}
max +------------+
    |   stack    |   <-- cresce verso il basso
    +------------+
    |     v      |
    |     ^      |
    +------------+
    |   heap     |   <-- cresce verso l'alto
    +------------+
    |   data     |
    +------------+
    |   text     |
  0 +------------+
\end{verbatim}

\dfn{Stack frame}{
  Ogni \textbf{funzione}, quando viene eseguita, ottiene un \textbf{frame} sullo stack.
  \begin{itemize}
    \item Frame \textbf{creato} quando viene chiamata la funzione, \textbf{distrutto} quando esce (exit/return).
    \item Contiene le \textbf{variabili locali}.
    \item Contiene le informazioni sul \textbf{flusso di controllo}: return address, old frame pointer.
    \item \textbf{Lo stack cresce verso il basso}.
  \end{itemize}
}

\subsection{Chiamata di una funzione}

Sia $P$ il chiamante (\textit{calling}) e $Q$ il chiamato (\textit{called}) da $P$.

\subsubsection{Lato chiamante (P)}

\begin{enumerate}
  \item Pusha i \textbf{parametri} per la funzione $Q$ sullo stack (tipicamente in \textbf{ordine inverso} di dichiarazione: \texttt{param 2}, poi \texttt{param 1}).
  \item Esegue l'istruzione \texttt{call} che invoca la funzione target $\Rightarrow$ pusha l'\textbf{indirizzo di ritorno} (\texttt{Return addr in P}) sullo stack.
\end{enumerate}

\subsubsection{Prologo della chiamata (Q)}

\begin{enumerate}
  \setcounter{enumi}{2}
  \item Pusha il \textbf{valore corrente del frame pointer} (che punta al frame del chiamante).
  \item Imposta il \textbf{frame pointer} al valore corrente dello stack pointer $\Rightarrow$ identifica la \textbf{nuova posizione del frame} per $Q$.
  \item \textbf{Alloca spazio per le variabili locali} muovendo lo stack pointer verso il basso.
  \item Esegue il \textbf{corpo} della funzione chiamata.
\end{enumerate}

\subsubsection{Layout del frame durante l'esecuzione di Q}

\begin{verbatim}
P:  +----------------------+
    |  Return addr         |
    |  Old frame pointer   |  <-----+
    |  param 2             |        |
    |  param 1             |        |
Q:  |  Return addr in P    |        |
    |  Old frame pointer   |  <-- Frame pointer
    |  local 1             |
    |  local 2             |  <-- Stack pointer
    +----------------------+
\end{verbatim}

\subsubsection{Uscita di Q}

\begin{enumerate}
  \setcounter{enumi}{6}
  \item Imposta lo \textbf{stack pointer} al valore del frame pointer (scartando lo spazio delle variabili locali).
  \item Pop dell'\textbf{old frame pointer} (ripristina il link al frame del chiamante).
  \item Esegue l'istruzione \texttt{ret}, che pop l'indirizzo salvato e restituisce il controllo a $P$: il \textbf{return address} \`e l'indirizzo dell'istruzione in $P$ immediatamente dopo la \texttt{call}.
\end{enumerate}

\subsubsection{Ripresa in P}

\begin{enumerate}
  \setcounter{enumi}{9}
  \item Pop dei \textbf{parametri} della funzione chiamata dallo stack.
  \item Continua l'esecuzione con l'istruzione successiva alla chiamata.
\end{enumerate}

\subsection{Esempio di Stack Frame}

\begin{verbatim}
0: string count(string s, character c) {
     integer counter;
     integer pos;
   ...
5: }
6: void main(integer argc, strings argv) {
7:   count("testing", "t");
8: }
\end{verbatim}

Lo stack frame durante l'esecuzione di \texttt{count}:
\begin{verbatim}
High +--------------+
     |   argv       |
     |   argc       |
     |   "t"        |   <-- main()
     |   "testing"  |
     |   IP = 8     |   <-- count(): tornare a riga 8
     |   counter    |
     |   pos        |
     +--------------+
Low                       Stack grows downward
\end{verbatim}

\nt{
  \textbf{Concetto chiave}: il frame deve tenere traccia dell'\textbf{IP successivo} (return address): la CPU lo user\`a a fine funzione per tornare al chiamante.
}

\subsection{Esempio di Ricorsione}

\begin{verbatim}
0: integer r(integer n) {
1:   if (n > 0) r(n - 1);
2:   return n;
3: }
4: void main(integer argc, strings argv) {
5:   r(3); // should return 3
6: }
\end{verbatim}

Stack risultante (dall'alto verso il basso):
\begin{verbatim}
main()  -> IP = ..., argv=3
r(3)    -> IP = 6,  n-1 = 2
r(2)    -> IP = 2,  n-1 = 1
r(1)    -> IP = 2,  n-1 = 0
r(0)    -> IP = 2,  0
\end{verbatim}

Ogni chiamata ricorsiva impila un nuovo frame.

\subsection{Review}

\begin{enumerate}
  \item I programmi in esecuzione sono nella \textbf{memoria (RAM)}.
  \item Il codice assembly \`e nella \textbf{program memory}. La CPU tiene traccia della prossima istruzione nel registro \textbf{IP / PC}.
  \item La \textbf{memoria dati} \`e strutturata come una \textbf{pila di frame}:
    \begin{itemize}
      \item Ogni invocazione di funzione (\texttt{call}) aggiunge un frame.
      \item Ogni frame contiene: \textbf{IP salvato} (e l'indirizzo del frame precedente) e le \textbf{variabili locali}.
    \end{itemize}
\end{enumerate}

\section{Stack Buffer Overflow}

\dfn{Stack Buffer Overflow}{
  Si verifica quando il \textbf{buffer} si trova nello \textbf{stack}, solitamente come variabile locale nello stack frame di una funzione. Detto anche \textbf{stack smashing}.

  Pi\`u formalmente: un processo (programma in esecuzione) tenta di memorizzare dati \textbf{oltre i limiti} di un buffer (sullo stack) di \textbf{dimensione fissa}, sovrascrivendo le posizioni di memoria adiacenti.
}

\nt{
  Visti per la prima volta nel \textbf{Morris Internet Worm (1988)}, sfruttando la funzione C \texttt{gets()} nel daemon \texttt{fingerd}.
}

\subsection{Un esempio semplice di programma vulnerabile}

\begin{verbatim}
0: void print(string s) {
     // only holds 32 characters, max
     string buffer[32];
1:   strcpy(buffer, s);   // copy s into buffer
2:   puts(buffer);        // print buffer to stdout
3: }

4: void main(integer argc, strings argv) {
5:   for (; argc > 0; argc = argc - 1) {
6:     print(argv[argc]);
7:   }
8: }
\end{verbatim}

Output normale:
\begin{verbatim}
$ ./print Hello World
World
Hello

$ ./print arg1 arg2 arg3
arg3
arg2
arg1
\end{verbatim}

\paragraph{Dov'\`e il problema?} \texttt{strcpy(buffer, s)} \textbf{non controlla la lunghezza} dell'input!
\begin{verbatim}
#include <cstring>
char *strcpy(char *to, const char *from);
\end{verbatim}
Da man-page: \textit{strcpy() does not perform bounds checking, and thus risks overrunning from or to. For a similar (and safer) function that includes bounds checking, see} \texttt{strncpy()}.

\subsection{Cosa succede se s \`e pi\`u lunga di 32 caratteri?}

Il buffer di 32 byte viene riempito, e i caratteri eccedenti \textbf{sovrascrivono}:
\begin{itemize}
  \item l'\textbf{old frame pointer}
  \item il \textbf{return address salvato} (\texttt{IP = 7})
  \item potenzialmente i parametri/dati di \texttt{main()}
\end{itemize}

\begin{verbatim}
High +--------------------+
     |   argv             |
     |   argc             |
     |   IP = ... (main)  |
     |   IP = 7 (print)   |  <-- sovrascritto!
     |   Data from argv   |  <-- buffer riempito oltre i limiti
Low  +--------------------+
\end{verbatim}

\subsection{Crash del programma e DoS}

Quando \texttt{print()} tenta di trasferire il controllo all'indirizzo di ritorno (\texttt{IP = 7}), salta a una \textbf{posizione di memoria illegale} $\Rightarrow$ \textbf{Segmentation Fault} $\Rightarrow$ terminazione anomala.

\nt{
  Fornire semplicemente un \textbf{input casuale} porta al crash del programma. Poich\'e il programma si \`e bloccato non pu\`o pi\`u fornire la funzione/servizio: nella sua forma pi\`u semplice, uno stack overflow \`e una forma di \textbf{Denial-of-Service (DoS)}.
}

\section{Smashing the Stack}

\subsection{Idea chiave}

\begin{itemize}
  \item I bug di buffer overflow possono \textbf{sovrascrivere gli IP salvati} $\Rightarrow$ crash del programma.
  \item \textbf{Idea chiave 1}: sostituire l'IP salvato $\Rightarrow$ pu\`o puntare ovunque l'attaccante voglia. \textbf{Ma dove?}
  \item \textbf{Idea chiave 2}: riempire il buffer con \textbf{codice dannoso}.
    \begin{itemize}
      \item Il codice macchina \`e solo una stringa di byte.
      \item Cambia l'IP per puntare al \textbf{codice dannoso (shellcode)} sullo stack.
    \end{itemize}
\end{itemize}

\subsection{Exploit v1 --- layout di memoria}

Indirizzi di memoria di esempio:
\begin{verbatim}
:1000  +--------------------+
       |   argv             |
:996   +--------------------+
       |   argc             |
:992   +--------------------+
       |   IP = ... (main)  |
:988   +--------------------+
       |   IP = 7 (print)   |  <-- da sovrascrivere
:984   +--------------------+
       |                    |
       |   buffer (32 B)    |
:952   +--------------------+
\end{verbatim}

\subsection{Exploit v1 --- sovrascrittura}

L'attaccante inietta:
\begin{itemize}
  \item nel buffer (a partire da \texttt{:952}): \textbf{shellcode} (codice malicious).
  \item dove si trova \texttt{IP = 7} (a \texttt{:984}): \texttt{IP = 952} $\Rightarrow$ punta all'inizio dello shellcode.
\end{itemize}

\begin{verbatim}
:1000  +--------------------+
       |   argv             |
:996   |   argc             |
:992   |   IP = ...         |
:988   |   IP = 952         |  <-- sovrascritto!
:984   +--------------------+
       |                    |
       | Malicious code /   |
       | Shellcode          |
:952   +--------------------+
\end{verbatim}

Quando \texttt{print()} ritorna, salta a \texttt{:952} ed esegue lo shellcode.

\section{Shellcode}

\dfn{Shellcode}{
  Componente essenziale di molti attacchi di buffer overflow: il \textbf{codice fornito dall'attaccante} a cui si trasferisce l'esecuzione, spesso salvato nel buffer in overflow.

  La sua funzione classica \`e quella di trasferire il controllo a un \textbf{interprete della riga di comando dell'utente, o shell}, dando accesso a qualsiasi programma disponibile sul sistema \textbf{con i privilegi del programma attaccato}.
}

\subsection{Caratteristiche}

\begin{itemize}
  \item Lo shellcode \`e semplicemente \textbf{machine code}: una serie di valori binari corrispondenti alle istruzioni macchina e ai dati che implementano la funzionalit\`a desiderata.
  \item \`E \textbf{specifico per architettura del processore} e \textbf{sistema operativo}, perch\'e deve interagire con le system call.
\end{itemize}

\paragraph{Esempio --- Shellcode minimale (in C):}
\begin{verbatim}
{
    // execute a shell with the privileges of the vulnerable program
    exec("/bin/sh", NULL);
}
\end{verbatim}

\subsection{Restrizioni generiche sul contenuto}

\thm{Vincoli sullo shellcode}{
  Uno shellcode pratico deve soddisfare tre vincoli:
  \begin{enumerate}
    \item \textbf{Position-Independent}: deve poter essere eseguito indipendentemente da dove si trovi in memoria. Possono essere usati solo riferimenti \textbf{relativi}, offset rispetto all'IP corrente.
    \item \textbf{No NULL bytes}: in C la stringa termina con NULL ($\Rightarrow$ \texttt{strcpy} si fermerebbe!). L'unico NULL ammesso \`e alla \textbf{fine}, dopo tutto il codice e l'indirizzo di ritorno sovrascritto.
    \item \textbf{Sopravvivenza alle trasformazioni}: cosa succede se il programma decifra/converte la stringa prima di copiarla? Se trasforma minuscole in maiuscole? Lo shellcode deve essere costruito per \textbf{evitare o tollerare} queste modifiche.
  \end{enumerate}
}

\subsection{Un classico shellcode per attacco execve}

Codice C equivalente:
\begin{verbatim}
int main(int argc, char *argv[]) {
    char *sh;
    char *args[2];
    sh = "/bin/sh";
    args[0] = sh;
    args[1] = NULL;
    execve(sh, args, NULL);
}
\end{verbatim}

Codice in assembly x86 (con commenti):
\begin{verbatim}
       nop
       nop
       jmp find        // jump to end of code
cont:  pop %esi        // pop address of sh off stack into %esi
       xor %eax, %eax  // zero contents of EAX
       mov %al, 0x7(%esi)   // copy zero byte to end of string sh
       lea (%esi), %ebx     // load address of sh into %ebx
       mov %ebx, 0x8(%esi)  // save address of sh in args[0]
       mov %eax, 0xc(%esi)  // copy zero to args[1]
       mov $0xb, %al        // copy execve syscall number (11) to AL
       mov %esi, %ebx       // copy address of sh into %ebx
       lea 0x8(%esi), %ecx  // copy address of args to %ecx
       lea 0xc(%esi), %edx  // copy address of args[1] to %edx
       int $0x80            // software interrupt to execute syscall
find:  call cont            // call cont -- saves next address on stack
sh:    .string "/bin/sh"
args:  .long 0
       .long 0
\end{verbatim}

\subsection{Generare shellcode}

Tool comuni:
\begin{itemize}
  \item \textbf{Metasploit Project} --- \url{https://www.metasploit.com/}
  \item \textbf{pwntools} --- \url{https://docs.pwntools.com/en/stable/}
\end{itemize}

\section{NOP Sled}

\subsection{Hitting the Target --- il problema}

\qs{Come possiamo indovinare in modo affidabile l'indirizzo dello shellcode?}{
  \begin{itemize}
    \item L'indirizzo dello shellcode deve essere \textbf{indovinato esattamente}: l'IP deve saltare all'inizio preciso dello shellcode.
    \item \textbf{Gli indirizzi dello stack cambiano spesso}: cambiano ogni volta che un programma viene eseguito.
  \end{itemize}
  \textbf{Trucco}: rendere il bersaglio \textbf{ancora pi\`u grande} cos\`i \`e pi\`u facile da colpire.
}

\subsection{NOP Sled}

\dfn{NOP Sled}{
  La maggior parte delle CPU supporta istruzioni \textbf{no-op}: istruzioni semplici, da un byte, che non fanno nulla. Su Intel x86 l'opcode \texttt{0x90} \`e il NOP.

  L'idea \`e di creare una \textit{slitta} di NOP davanti allo shellcode: funge da grande \textbf{rampa}. Se l'IP atterra in un punto \textbf{qualsiasi} della rampa, esegue NOP finch\'e non raggiunge lo shellcode.
}

\subsection{Exploit v2 con NOP sled}

Buffer pi\`u grande (128 byte) consente:
\begin{verbatim}
:1000  +--------------------+
       |   argv             |
:996   |   argc             |
:992   |   IP = ...         |
:988   +--------------------+
:984   |   IP = 900         |  <-- punta nel mezzo del NOP sled
       +--------------------+
       |   Malicious code   |
       |   (shellcode)      |
       +--------------------+
       |                    |
       |   NOP sled         |
       |   (0x90 0x90 ...)  |
:856   +--------------------+
\end{verbatim}

Anche se l'attaccante "indovina male" e atterra in un punto qualsiasi del NOP sled, l'esecuzione \textit{scivola} fino allo shellcode.

\section{Funzioni C non sicure e altre forme di overflow}

\subsection{Routine comuni della libreria C non sicure}

\begin{center}
\begin{tabular}{|l|p{8cm}|}
\hline
\textbf{Funzione} & \textbf{Descrizione} \\
\hline
\texttt{gets(char *str)} & Read line from standard input into \texttt{str} \\
\texttt{sprintf(char *str, char *fmt, ...)} & Create \texttt{str} according to supplied format and variables \\
\texttt{strcat(char *dest, char *src)} & Append contents of \texttt{src} to \texttt{dest} \\
\texttt{strcpy(char *dest, char *src)} & Copy contents of \texttt{src} to \texttt{dest} \\
\texttt{vsprintf(char *str, char *fmt, va\_list ap)} & Create \texttt{str} according to format and variables \\
\hline
\end{tabular}
\end{center}

\nt{
  Tutte queste \textbf{non controllano i limiti} del buffer di destinazione.
}

\subsection{Altre forme di attacchi di overflow}

\begin{itemize}
  \item \textbf{Heap Overflows}: overflow in memoria allocata dinamicamente.
  \item \textbf{Global Data Area Overflows}: overflow in variabili globali.
  \item \textbf{Format String Overflows}: sfruttamento di funzioni \texttt{printf}-like con format string controllata dall'utente.
  \item \textbf{Integer Overflows}: overflow numerico che produce dimensioni errate.
\end{itemize}

\section{Difese contro i Buffer Overflow}

\subsection{Stack Canaries --- StackGuard}

\dfn{Stack Canary}{
  Il compilatore aggiunge un \textbf{valore speciale} (sentinel) sullo stack \textbf{prima di ogni IP salvato}.
  \begin{itemize}
    \item Il canary \`e impostato a un \textbf{valore casuale} in ogni frame.
    \item All'\textbf{uscita della funzione} il canary viene \textbf{controllato}.
    \item Se \`e stato alterato, il programma esce con un \textbf{errore}.
  \end{itemize}
}

Esempio:
\begin{verbatim}
0: void print(string s) {
1:   __set_stack_canary(random());
2:   string buffer[32];
3:   strcpy(buffer, s);
4:   puts(buffer);
5:   __check_stack_canary()
6: }
\end{verbatim}

Layout di memoria con canary:
\begin{verbatim}
:1000  +--------------------+
       |   argv             |
:996   |   argc             |
:992   |   IP = ...         |
:988   +--------------------+
       |   IP = 7           |
:980   +--------------------+
       |   canary = 189476  |  <-- sentinel
       +--------------------+
       |                    |
       |   buffer           |
:948   +--------------------+
\end{verbatim}

\mprop{Effetto del canary contro lo smashing}{
  Il malicious code sovrascrive il canary insieme all'IP. Al \texttt{\_\_check\_stack\_canary()} la funzione rileva la modifica $\Rightarrow$ chiama \texttt{exit()} invece di tornare $\Rightarrow$ \textbf{l'attacco fallisce}.
}

\subsection{Stack non eseguibile}

\dfn{NX / DEP (Data Execution Prevention)}{
  Le CPU moderne impostano la memoria dello stack come \textbf{lettura/scrittura}, \textbf{ma non eXecute}.
}

\nt{
  \textbf{Non evita} i buffer overflow ma \textbf{impedisce che lo shellcode venga inserito ed eseguito nello stack}.
}

\subsection{Address-Space Layout Randomization (ASLR)}

\dfn{ASLR}{
  Difesa al \textbf{livello del sistema operativo}: \textbf{randomizzazione del layout della memoria}. Ogni volta che un programma viene caricato, la posizione del codice (program memory) e dei dati (data memory) cambia.
}

\begin{itemize}
  \item Rende \textbf{pi\`u difficile} per l'attaccante indovinare la destinazione del buffer sullo stack.
  \item \textbf{Non impedisce lo sfruttamento, ma lo rende pi\`u difficile}.
  \item Supportato da tutti i moderni sistemi operativi.
  \item Funziona meglio quando la dimensione della memoria \`e grande (pi\`u entropia $\Rightarrow$ pi\`u difficile indovinare).
\end{itemize}

\subsection{Write Safe Code!}

\nt{
  Il messaggio importante \`e che \textbf{se i programmi non sono codificati correttamente} per proteggere le loro strutture dati fin dall'inizio, sono soggetti a potenziali attacchi.

  Sebbene canary, NX e ASLR possano bloccare molti attacchi, alcuni --- come la corruzione del valore di una variabile adiacente per alterare il comportamento del programma --- \textbf{non possono essere bloccati se non con codice sicuro}.
}

Linee guida:
\begin{itemize}
  \item Usare funzioni con bounds checking: \texttt{strncpy}, \texttt{snprintf}, \texttt{fgets} invece di \texttt{strcpy}, \texttt{sprintf}, \texttt{gets}.
  \item Validare gli input.
  \item Evitare linguaggi/funzioni che lasciano la responsabilit\`a della memoria al programmatore quando possibile (es. usare Rust).
\end{itemize}

% \end{document}
